\scanpage{52}
\mspage{8}
We thus have the notion of a \uline{weighted graph} \deleted{\unsure{enriched}} \inserted{\uline{with markings}} $\Gamma$

\diagram{\begin{tikzcd}[ampersand replacement=\&,column sep=large,row sep=large]
\vec R_l^+\ \deleted{\unsure{\mathbb D_\infty|D_\infty}}\subset\underset{\sigma_0,\sigma_1}{\vec R}
 \arrow[r]\&\vec A\arrow[d,"\unsure{o|\sigma}"']\arrow[loop right,"\sigma"]\\
\&S\arrow[r,"\gamma"]\&\mathbb N
\end{tikzcd}}

Groupoid (\uline{Gr.pndrep})

functor
\[
 S_*\mathcal TD_{\mathrm{top}}^\pm\longrightarrow(\mathrm{Grpondrep})
\]
which will probably be important in making the operation of ``gluing'' explicit.
It is clear in any case that any algebro-combinatorial description-translation of
$S\mathcal TD_{\mathrm{top}}^\pm$ will have to make it possible, for each object of
\deleted{the category} \inserted{groupoid-}paradigm, to reconstruct (up to unique isomorphisms) the corresponding \uline{weighted graph with markings}.
\inserted{--- more precisely, one must reconstruct the canonical functor
$S_*\mathcal TD_{\mathrm{top}}^\pm\longrightarrow(\mathrm{Gr.pond.rep.})$
(and not only the functor \unsure{with|:} values in $(\mathrm{Grbipond})$)}

\uline{Essential operations}

A) Operations of ``forgetting'' or of ``\uline{erasure}''.

\textcircled{1} \uline{Erasure of the markings} To an object $(X,\Sigma,R)$, equipped in addition with a subset $R'\subset R$, it associates $(X,\Sigma,R\setminus R')$ (erasure of $R'$).
\mspage{51}
